\chapter{De Rham Cohomology}

All sections and chapters above are just MATH5020, a graduate calculus course. Now we will switch to topology, starting MATH5070.

To motivate, recall in \(\R^3\), a vector field \(\vec{F}\) is said to be \begin{itemize}
    \item \emph{conservative} if \(\vec{F} = \nabla f\) for some \emph{potential function} \(f\).
    \item \emph{curl-free} if \(\nabla\times\vec{F} = 0\).
\end{itemize}
Clearly, a conservative vector field is curl-free by taking curl on both sides. But the converse is not always true.

\begin{example}
    Consider \[
        \vec{F} = \Big\<\frac{-y}{x^2 + y^2}, \frac{x}{x^2 + y^2}\Big\>
    \] defined on \(\R^2\setminus\{0\}\), it is not a gradient field but is curl-free.
\end{example}

In general, such converse is true if the domain of definition is \emph{simply-connected}. This is quite mysterious when we first encountered, but it will be precise once we learn the language of de Rham cohomology.

\section{De Rham Cohomology}

Let \(M\) be a smooth manifold, we have a map \[
    d:\Omg^k(M)\to\Omg^{k+1}(M)
\] where \(d^2 = 0\).
Then we obtain the \emph{de Rham complex} \[
    \begin{tikzcd}
        0 \ar{r} &\Omg^0(M)
        \ar{r}{d} &\Omg^1(M)
        \ar{r}{d} &\cdots
        \ar{r}{d} &\Omg^n(M) \ar{r} &0.
    \end{tikzcd}
\]

For manifolds, we can generalize the said notions as follows.
\begin{definition}
    A differential \(k\)-form \(\omega\) on a smooth manifold \(M\) is said to be \begin{enumerate}
        \item \emph{closed} if \(d\omega = 0\),
        \item \emph{exact} if there exists a \((k-1)\)-form \(\eta\) such that \(\omega = d\eta\).
    \end{enumerate}
    Denote by the vector spaces \[
        Z^k(M) := \ker\Big(d:\Omg^k(M)\longrightarrow\Omg^{k+1}(M)\Big),\quad
        B^k(M) := \im\Big(d:\Omg^{k-1}(M)\longrightarrow\Omg^k(M)\Big),
    \] which are the spaces of closed \(k\)-forms and exact \(k\)-forms, respectively. Note that \(B^k(M)\subset Z^k(M)\) since \(d^2 = 0\).
\end{definition}

These two spaces are not always the same, if we want to measure the difference between them, we take their quotient space, which is the de Rham cohomology group.
\begin{definition}[\emph{de Rham cohomology group}]
    The quotient vector space \[
        H_{\dR}^k(M) := Z^k(M) \Big/ B^k(M)
    \] is called the \(k\)-th \emph{de Rham cohomology group} of \(M\).
\end{definition}
Note that if \(k > n\), then \(B^k(M) = Z^k(M) = \{0\}\), so \(H_{\dR}^k(M) = \{0\}\).

\begin{definition}[\emph{Betti number}, \emph{Euler characteristic}]
    If \(\dim(H_{dR}^k(M)) <\infty\), then we define \[
        b_k(M) := \dim H_{\dR}^k(M)
    \] to be the \(k\)-th \emph{Betti number} of \(M\). Furthermore, we define the \emph{Euler characteristic} of \(M\) to be \[
        \chi(M) := \sum_{k=1}^n(-1)^k b_k(M).
    \]
\end{definition}
These two are topological invariants. 

\subsection*{The Ring Structure of Cohomology}

\begin{definition}
    Given \(\omg\in Z^k(M)\), we call \([\omg] \in H_{dR}^k(M)\) the \emph{cohomology class} of \(\omg\).
\end{definition}

\begin{definition}[\(\wed\) on \(H^k(M)\)]
    Let \(\omg\in Z^k(M)\) and \(\eta\in Z^l(M)\) so that \([\omg]\in H^k(M)\) and \([\eta]\in H^l(M)\), we define \[
        [\omg]\wed[\eta] = [\omg\wed\eta]\in H^{k+l}(M).
    \]
\end{definition}
\begin{proof}
    \textbf{Well-defindedness.}
    \begin{align*}
        d(\omg\wed\eta) = \underbrace{d\omg}_0\wed\eta + (-1)^k\omg\wed \underbrace{d\eta}_0 = 0,
    \end{align*} so \(\omg\wed\eta\) is closed.
    Let \([\omg_1] = [\omg_2]\in H^k(M)\), then \(\omg_2 = \omg_1 + d\xi\), as \begin{align*}
        \omg_2\wed\eta
        = (\omg_1 + d\xi)\wed\eta
        = \omg_1\wed\eta + d\xi\wed\eta,
    \end{align*} we need to show that \(d\xi\wed\eta\) is exact, which is again by \begin{align*}
        d(\xi\wed\eta)
        = d\xi\wed\eta + (-1)^k\xi\wed \underbrace{d\eta}_0
        = d\xi\wed\eta.
    \end{align*}
\end{proof}

Now, this wedge product \(H^k(M)\times H^l(M)\xlongrightarrow{\wed} H^{k+l}(M)\) makes \[
    H^*(M) := \bigoplus_{k=0}^n H^k(M)
\] a \emph{graded ring}.

Suppose \(\f:M\longrightarrow N\) is a smooth map, recall that \(\f^*\circ d = d\circ\f^*\), meaning that \[
    \f^*(Z^k(N))\subset Z^k(M),\quad \f^*(B^k(N))\subset B^k(M),
\]

This induces a map \(\f^*:H^k(N)\longrightarrow H^k(M)\) by \([\omg]\longmapsto [\f^*\omg]\)

\begin{lemma}\
    \begin{enumerate}
        \item \((\psi\circ\f)^* = \f^*\circ\psi^*\),
        \item \(\id^* = \id\).
    \end{enumerate}
\end{lemma}

\begin{theorem}
    If \(\f:M\to N\) is a diffeomorphism, then \[
        \f^*:H^k(N)\xrightarrow{\cong} H^k(M)
    \] is a linear isomorphism.
\end{theorem}

\subsection*{Computation of Cohomology Groups}

The remaining of this section is to compute these invariants for some manifolds and see when the \(\dim(H_{\dR}^k(M))\) is finite.

\begin{proposition}\label{prop:H0_conn_comp}
    Let \(M\) be a smooth manifold with \(r\) connected components, then \(H^0(M)\cong \R^r\).
\end{proposition}
\begin{proof}
    Since \(B^0(M) = \{0\}\), we have \begin{align*}
        H_{dR}^0(M) &= Z^0(M) \Big/ B^0(M) \\
        &\cong Z^0(M) \subset \Omg^0(M) = C^\infty(M).
    \end{align*}
    Note that
    \[
        0 = df = \sum \frac{\pa f}{\pa x^i}dx^i
        \implies \frac{\pa f}{\pa x^i} = 0 \text{ for all } i,
    \]
    which means \(f\) is locally constant. Since \(M\) has \(r\) connected components, \(f\) can take different constant values on each component, so we can represent \(f\) by a vector in \(\R^r\), each coordinate corresponding to the value of \(f\) on each component. This gives an isomorphism \(Z^0(M)\cong \R^r\), and hence \(H^0(M)\cong \R^r\).
\end{proof}

\begin{proposition}\label{prop:Hk_vanish}
    Let \(M^m\) be a smooth manifold, then \(H^k(M) = 0\) for \(k > m\).
\end{proposition}
\begin{proof}
    For each \(p\in M\), \(\dim T_pM = m\). If \(\omg\) is a \(k\)-form on \(M\), then \(\omg_p\in\bigwedge^kT_pM\), it suffices to show \(\bigwedge^kT_pM = 0\) for \(k > m\). Let \(\{e_1,\dots,e_m\}\) be a basis of \(T_pM\), then \[
        \{e_{i_1}\wed\dots\wed e_{i_k} : 1\le i_1 < \cdots < i_k \le m\}
    \] is a basis of \(\bigwedge^kT_pM\). However, if \(k > m\), at least two of the indices \(i_1,\dots,i_k\) must be equal, so \(e_{i_1}\wed\dots\wed e_{i_k} = 0\). Hence \(\bigwedge^kT_pM = 0\) for \(k > m\), and thus \(H^k(M) = 0\) for \(k > m\).
\end{proof}

\begin{example}
    For \(M = \R\), we draw the de Rham complex \[
        0 \longrightarrow \Omg^0(\R) \xlongrightarrow{d} \Omg^1(\R) \longrightarrow 0.
    \]
    Since the real line \(\R^1\) is connected, by Proposition~\ref{prop:H0_conn_comp}, we have \(H^0(\R^1)\cong \R\).  It remains to compute \(H^1(\R^1)\). 
    \begin{roof}
        We already know \(Z^1(M) = \Omg^1(\R)\). Let \(\omg = g dx\in\Omg^1(\R)\), let \[
            G(x) = \int_0^x g(t)dt,
        \] then \[
            dG = g dx = \omg,
        \]
        Hence \(B^1(M) = \Omg^1(\R)\), which forces \(H^1(M) = 0\).
    \end{roof}
    In summary, we have \[
        H^k(\R^1) \cong \begin{cases}
            \R & k = 0,\\
            0 & k > 0.
        \end{cases}
    \]
\end{example}

\newpage

\begin{example}
    Let \(S^1\) be the unit circle embedded in \(\R^2\). Since \(S^1\) is connected, \(H^0(S^1)\cong \R\). By Proposition~\ref{prop:Hk_vanish}, we have \(H^k(S^1) = 0\) for \(k > 1\). It remains to compute \(H^1(S^1)\).
    \begin{roof}
        We know \begin{align*}
            Z^1(S^1) &= \Omg^1(S^1) = \Big\{g\,d\t\Bigm| g\in C^\infty(S^1)\Big\}, \\
            B^1(S^1) &= \Big\{df\Bigm| f\in C^\infty(S^1)\Big\}.
        \end{align*}
        We claim that \[
            g\,d\t = df \iff \int_0^{2\pi} g(\t)d\t = 0.
        \]
        \begin{roof}
            \((\Rightarrow)\) is just the Stoke's theorem, i.e. \[
                \int_0^{2\pi} g(\t)\,d\t = \int_{S^1} df = 0.
            \]
            For \((\Leftarrow)\), let \(f(\t) = \int_0^\t g(t)\,dt\) be a function on \(\R\), by assumption, it is \(2\pi\)-periodic, therefore \[
                df = g\,d\t.
            \]
            Thus, \[
                B^1(S^1) = \Big\{
                    g\,d\t \Bigm|
                    g\in C^\infty(\R), g(x+2\pi) = g(x), \int_0^{2\pi} g(\t)d\t = 0   
                \Big\}
            \]
            Thus, we claim that \[
                \f = \int\cdot : H^1(S^1) \longrightarrow \R, \quad [g\,d\t] \longmapsto \int_0^{2\pi} g(\t)d\t
            \] is an isomorphism. Since \(\ker\f = 0\), it is injective. For surjectivity, given \(c\in\R\), we can take \(g(\t) = \frac{c}{2\pi}\) to get \(\f([g\,d\t]) = c\). Hence \(\f\) is surjective. Thus, \(H^1(S^1) = \R\). (Remark: \(d\t\) is not exact.)
        \end{roof}
        % \[
        %     g(\tha)d\tha = df
        %     \iff \int_0^{2\pi} g(\tha)d\tha = \int_0^{2\pi} df = f(2\pi) - f(0) = 0,
        % \] where \(f(\tha) = \int_0^\tha g(t)dt\).
        % Hence \(\f = \int_0^{2\pi}\cdot\) gives an isomorphism \(H^1(S^1) \to \R^1\). In summary, we have \[
        %     H^k(S^1) \cong \begin{cases}
        %         \R & k \leq 1,\\
        %         0 & k > 1.
        %     \end{cases}
        % \]
    \end{roof}
\end{example}




\section[Homotopy Invariance]{Homotopy Invariance\headingnote{Madsen, Tornhave, \emph{From Calculus to Cohomology}, \S6}}

\newcommand{\simH}{\underset{\mathrm{hom}}{\sim}}

\subsection*{Homotopy of Maps and Homotopy Equivalence}

\begin{definition}[\emph{homotopic maps}]
    Let \(f:M\longrightarrow N\) and \(g:M\longrightarrow B\) are \emph{homotopic} if there exists a smooth map \[
        \Phi:M\times[0,1]\longrightarrow N, \quad (p,t)\longmapsto \Phi(p,t)
    \] such that \[
        \Phi(p,0) = f(p),\quad \Phi(p,1) = g(p),\quad \forall p\in M,
    \]
    called the \emph{homotopgy} between \(f\) and \(g\). In this case, we write \(f\simH g\).
    % local perturbation of each others
\end{definition}

\begin{example}[\emph{straight-line homotopy}]
    Let \(f,g: M\longrightarrow\R^n\) be two smooth maps, define \(\Phi:M\times\R\longrightarrow\R^n\) by \[
        \Phi(x,t) = (1-t)f(x) + tg(x).
    \] Then \(\Phi\) is a homotopy between \(f\) and \(g\), called the \emph{straight-line homotopy}.
\end{example}

\begin{definition}[\emph{homotopy equivalence}]
    Two smooth manifolds \(M\) and \(N\) are \emph{homotopy equivalent} if there exist smooth maps \[
        f:M\longrightarrow N,\quad g:N\longrightarrow M
    \] such that \[
        g\circ f \simH \id_M,\quad f\circ g \simH \id_N.
    \]
\end{definition}

\begin{example}
    \(S^{n-1} \simH \R^n\setminus\{0\}\).
    \begin{roof}
        Consider the smooth maps \begin{alignat*}{2}
            \f=\iota: S^{n-1}
                &\longrightarrow \mathbb{R}^n\setminus\{0\},
            &\qquad \psi: \mathbb{R}^n\setminus\{0\}
                &\longrightarrow S^{n-1},\\
            x &\longmapsto x,
            & x &\longmapsto \frac{x}{\|x\|}.
        \end{alignat*}
        We have \(\psi\circ\f = \id\) because \[
            \begin{tikzcd}
                x \ar[|->]{r}{\f} & x \ar[|->]{r}{\psi} & \frac{x}{1} = x.
            \end{tikzcd}
        \]
        To show \(\f\circ\psi \simH \id\), we use the straight-line homotopy \[
            \Phi(x,t) = (1 - t)\frac{x}{\|x\|} + tx
            = \left(\frac{1 - t}{\|x\|} + t\right)x.
        \] Since the coefficient of \(x\) is positive on \(t\in [0,1]\), we see \(\Phi\) is well-defined into \(\R^n\setminus\{0\}\).
    \end{roof}
\end{example}

\begin{example}
    A \emph{star-shaped} region \(U\) in \(\R^n\) is a region with a \emph{center} \(x_0\in U\) satisfying \[
        (1-t)x + tx_0 \in U \quad \forall x\in U,\ \forall t\in[0,1].
    \]
    We claim \(U\simH\{x_0\}\). Consider the inclusion and restriction maps \[
        \f:\{x_0\}\longrightarrow U, \quad \psi:U\longrightarrow\{x_0\}.
    \]
    Clearly \(\psi\circ\f = \id:\{x_0\} \longrightarrow \{x_0\}\), and \[
        \f\circ\psi : U\longrightarrow U, \quad x\longmapsto x_0.
    \] is homotopic to \(\id:U\longrightarrow U\) by \[
        \Phi(x,t) = (1-t)x + tx_0 ,
    \] which is well-defined into \(U\) by its star shape.
\end{example}

We see that a star-shaped region is homotopy equivalent to a point. Actually, we have a name for this kind of space.

\begin{definition}
    Any region homotopic to a single-point set is called \emph{contractible}.
\end{definition}

\subsection*{Homotopy Invariance of Cohomology}
\begin{theorem}
    Let \(f,g:M\longrightarrow N\) be homotopic, then \[
        f^* = g^*:H^k(N)\longrightarrow H^k(M).
    \]
\end{theorem}
\begin{proof}
    \begin{roof}
        Consider the diagram of de Rham cochain complexes \[
            \begin{tikzcd}[column sep=8em]
                \vdots \ar{d}{d} & \vdots \ar{d}{d} \\
                \Omg^{k-1}(N) \ar{r} \ar{d}{d} & \Omg^{k-1}(M) \ar[blue]{d}{d} \\
                \Omg^k(N) \ar{r}{f^*}[']{g^*} \ar[blue]{ur}{\wt{Q}} \ar[red]{d}{d} & \Omg^k(M) \ar{d}{d} \\
                \Omg^{k+1}(N) \ar{r} \ar[red]{ur}[']{\wt{Q}} \ar{d}{d} & \Omg^{k+1}(M) \ar{d}{d} \\
                \vdots & \vdots
            \end{tikzcd}
        \]
        The strategy is to construct \(\wt{Q}:\Omg^k(N)\longrightarrow\Omg^{k-1}(M)\) such that \(g^* - f^* = \textcolor{blue}{d\wt{Q}} + \textcolor{red}{\wt{Q}d}\).
        Afterwards, let \(\omg\in Z^k(N)\), then \begin{align*}
            g^*\omg - f^*\omg
            = d\wt{Q}(\omg) + \wt{Q}(\cancel{d\omg}) = d\wt{Q}(\omg)\in B^k(M),
        \end{align*} which implies \((g^* - f^*)([\omg]) = 0\in H^k(M)\).
    \end{roof}
    We first do a special case.
    \begin{roof}
        Let \(X\) be a complete vector field on \(M\) and \(\phi_t\) the flow generated by \(X\). We want to construct a linear operator \(
            Q:\Omg^k(M)\longrightarrow\Omg^{k-1}(M)
        \) such that \[
            \phi_1^* - \phi_0^* = dQ + Qd.
        \]
        \begin{roof}
            We apply the fundamental theorem of calculus, \begin{align*}
                \phi_1^*\omg - \phi_0^*\omg
                = \int_0^1 \frac{d}{dt}\phi_t^*\omg\,dt.
            \end{align*}
            Then by the Cartan magic formula~\ref{thm:cartan_magic_formula}, \begin{align*}
                \frac{d}{dt} \phi_t^*\omg
                = \frac{d}{ds} \phi_{t+s}^*\omg \bigg|_{s=0}
                = \frac{d}{ds} \phi_s^*\phi_t^*\omg \bigg|_{s=0}
                &= \L_X \phi_t^*\omg \\
                &= d(\iota_X\phi_t^*\omg) + \iota_X(d\phi_t^*\omg) \\
                &= d\underbrace{(\iota_X\phi_t^*\omg)}_{Q_t(\omg)} + \underbrace{(\iota_X\phi_t^*)}_{Q_t}\,d\omg.
            \end{align*}
            Put \(Q_t = \i_X\phi_t^*\), then \(\dfrac{d}{dt}\phi_t^*\omg = dQ_t(\omg) + Q_t(d\omg)\), finally, define \[
                \Omg^k(M) \xlongrightarrow{Q} \Omg^{k-1}(M),\quad \omg\longmapsto \int_0^1 Q_t(\omg)\,dt,
            \] then indeed \begin{align*}
                \phi_1^*\omg - \phi_0^*\omg
                &= \int_0^1 dQ_t(\omg)\,dt + \int_0^1 Q_t(d\omg)\,dt \\
                &= d\int_0^1 Q_t(\omg)\,dt + \int_0^1 Q_t(d\omg)\,dt
                = dQ(\omg) + Q(d\omg).
            \end{align*}
        \end{roof}
    \end{roof}
    We do the general case.
    \begin{roof}
        Put \(W = M\times\R\) and \(f,g\) be homotopic by \(\Phi:M\times[0,1]\longrightarrow N\) such that \[
            \Phi(p,0) = f(p),\quad \Phi(p,1) = g(p)\qquad \forall p\in M.
        \]
        Note that \(X = \frac{\pa}{\pa t}\) is a complete vector field on \(W\) whose flow \(\phi_t(p,a) = (p, a + t)\). Denote by \[
            \i:M\longrightarrow W,\quad \i(p) = (p,0)
        \] the inclusion map, then \(f = \Phi\circ\phi_0\circ\i\) and \(g = \Phi\circ\phi_1\circ\i\). Applying the special case, we obtain \begin{align*}
            g^*\omg - f^*\omg
            &= \i^*\phi_1^*\Phi^*\omg - \i^*\phi_0^*\Phi^*\omg \\
            &= \i^*(\phi_1^* - \phi_0^*)\Phi^*\omg \\
            &= \i^*(dQ + Qd)\Phi^*\omg \\
            &= (d\wt{Q} + \wt{Q}d)\omg,
        \end{align*} where we put \(\wt{Q} := \i^*\circ Q\circ\Phi^*\).
    \end{roof}
\end{proof}

This results in the main result of the section, the de Rham cohomology is a \emph{homotopy invariant}, and hence a \emph{topological invariant}, even when the spaces are not diffeomorphic.
\begin{corollary}
    If \(M\) and \(N\) are homotopy equivalent, then \[
        H^k(M)\cong H^k(N).
    \]
\end{corollary}
\begin{proof}
    By definition, there exist \(\f,\psi\) such that \(\f\circ\psi \simH \id_N\) and \(\psi\circ\f \simH \id_M\). By the theorem, we have \[
        (\f\circ\psi)^* = (\id_N)^* = \id_{H^k(N)},\quad
        (\psi\circ\f)^* = (\id_M)^* = \id_{H^k(M)}.
    \] Hence, \(\f^*:H^k(N)\longrightarrow H^k(M)\) and \(\psi^*:H^k(M)\longrightarrow H^k(N)\) are inverse to each other, which gives the desired isomorphism.
\end{proof}

\begin{corollary}
    If \(M\) is contractible, then \[
        H^k(M) = \begin{cases}
            \R & k = 0,\\
            0 & k > 0.
        \end{cases}
    \]
\end{corollary}
\begin{proof}
    Simply note that \(M\simH\{p\}\simH \R\).
\end{proof}

% \begin{definition}
%     \[
%         \omg - \omg' = d\tau\in B^k(M),
%     \] i.e.\ differ by an exact form. In this case, we say \(\omg\) and \(\omg'\) are \emph{cohomologous}.
% \end{definition}

\section{Mayer-Vietoris Sequence}

\subsection*{Cohomological Algebra\headingnote{Tu, \emph{An Introduction to Manifolds}, \S25}}

Now, let's temporary forget about those differential forms and exterior derivative, and look at the most basic setup.

\begin{definition}
    A \emph{co-chain complex} is a sequence of vector spaces \[
        \begin{tikzcd}
            \cdots \ar{r} & C^{k-1} \ar{r}{d_{k-1}} & C^k \ar{r}{d_k} & C^{k+1} \ar{r}{d_{k+1}} & \cdots
        \end{tikzcd}
    \] where \(d\) is a linear map such that \(d^2 = 0\). Furthermore, a complex is \emph{exact} if \(\im d_k = \ker d_{k+1}\).
\end{definition}

Here are two ingredients for diagram chasing arguments.

\begin{definition}
    A sequence \[
        \begin{tikzcd}
            A \ar{r}{f} & B \ar{r}{g} & C
        \end{tikzcd}
    \] is said to be \emph{exact at \(B\)} if \(\im f = \ker g\). A sequence is \emph{exact} if it is exact at every term except the first and last. A \emph{short} sequence is a sequence of the form \[
        \begin{tikzcd}
            0 \ar{r} & A \ar{r}{f} & B \ar{r}{g} & C \ar{r} & 0.
        \end{tikzcd}
    \]
\end{definition}

We see that a short exact sequence already contains a lot of information, for \[
    \begin{tikzcd}
        0 \ar{r}{\i} & A \ar{r}{f} & B \ar{r}{g} & C \ar{r}{\pi} & 0,
    \end{tikzcd}
\] we see that since \(\dom\i = 0\) and \(\im\pi = 0\), we have \(\i = \id\) and \(\pi = 0\), so we can omit them. Furthermore, \begin{itemize}
    \item exactness at \(A\) means \(\im\i = 0 = \ker f\), hence \(f\) is injective,
    \item exactness at \(C\) means \(\im g = C = \ker\pi\), hence \(g\) is surjective.
\end{itemize}
Therefore, we can also write the short exact sequence as \[
    \begin{tikzcd}
        0 \ar{r} & A \ar[hook]{r}{f} & B \ar[->>]{r}{g} & C \ar{r} & 0.
    \end{tikzcd}
\]

\begin{definition}
    A \emph{commutative diagram} is a diagram \[
        \begin{tikzcd}
            C_1 \ar{r}{f} \ar{d}{g_1} & C_2 \ar{d}{g_2} \\
            C_1' \ar{r}{f'} & C_2'
        \end{tikzcd}
    \]
\end{definition} such that \(g_2\circ f_1 = h_1\circ g_1\). In other words, the two paths from \(C_1\) to \(C_2'\) are equal.

\begin{lemma}
    Suppose we have a commutative diagram of short exact sequences \[
        \begin{tikzcd}
            0 \ar{r} & A_1 \ar[hook]{r}{f_1} \ar{d}{g_1} & B_1 \ar[->>]{r}{h_1} \ar{d}{g_2} & C_1 \ar{r} \ar{d}{g_3} & 0 \\
            0 \ar{r} & A_2 \ar[hook]{r}{f_2} & B_2 \ar[->>]{r}{h_2} & C_2 \ar{r} & 0
        \end{tikzcd}
    \]
\end{lemma}

Short sequences of co-chain complexes can induce a long exact sequence in cohomology. Suppose we have a family of short exact sequences \[
    \begin{tikzcd}
         & \vdots & \vdots & \vdots &  \\
        0 \ar{r} & C_1^{k+1} \ar{r}{i} & C_2^{k+1} \ar{r}{j} & C_3^{k+1} \ar{r} & 0 \\
        0 \ar{r} & C_1^k \ar{r}{i} & C_2^k \ar{r}{j} & C_3^k \ar{r} & 0 \\
        0 \ar{r} & C_1^{k-1} \ar{r}{i} & C_2^{k-1} \ar{r}{j} & C_3^{k-1} \ar{r} & 0 \\
         & \vdots & \vdots & \vdots & 
    \end{tikzcd}
\]
We can connect them vertically with \(C^k\xrightarrow[\hspace{1.5em}]{d} C^{k+1}\) to obtain a commutative digram \[
    \begin{tikzcd}
         & \vdots & \vdots & \vdots & \\
        0 \ar{r} & C_1^{k+1} \ar{r}{i} \ar[blue]{u}{d} & C_2^{k+1} \ar{r}{j} \ar[blue]{u}{d} & C_3^{k+1} \ar{r} \ar[blue]{u}{d} & 0 \\
        0 \ar{r} & C_1^k \ar{r}{i} \ar[blue]{u}{d} & C_2^k \ar{r}{j} \ar[blue]{u}{d} & C_3^k \ar{r} \ar[blue]{u}{d} & 0 \\
        0 \ar{r} & C_1^{k-1} \ar{r}{i} \ar[blue]{u}{d} & C_2^{k-1} \ar{r}{j} \ar[blue]{u}{d} & C_3^{k-1} \ar{r} \ar[blue]{u}{d} & 0 \\
        & \vdots \ar[blue]{u}{d} & \vdots \ar[blue]{u}{d} & \vdots \ar[blue]{u}{d} &
    \end{tikzcd}
\] where each column is a co-chain complex. By diagram chasing, we can construct the \emph{boundary maps} \(C_3^k\xrightarrow[\hspace{1.5em}]{\delta} C_1^{k+1}\) to obtain

\tikzcdset{
  boundary/.style={
    green!55!black,
    out=80, in=-80, looseness=.5,
    start anchor=north east,
    end anchor=south east,
    preaction={draw=white, line width=3pt, arrows=-}
  }
}

\[
    \begin{tikzcd}[row sep=3em, column sep=3em]
    & \vdots & \vdots & \vdots & \\
    0 \arrow[r]
        & C_1^{k+1} \ar{r}{i} \ar[blue]{u}{d}
        & C_2^{k+1} \ar{r}{j} \ar[blue]{u}{d}
        & C_3^{k+1} \ar{r} \ar[blue]{u}{d}
        & 0 \\
    0 \arrow[r]
        & C_1^k \ar{r}{i} \ar[blue]{u}{d}
        & C_2^k \ar{r}{j} \ar[blue]{u}{d}
        & C_3^k \ar{r} \ar[blue]{u}{d}
        & 0 \\
    0 \arrow[r]
        & C_1^{k-1} \ar{r}{i} \ar[blue]{u}{d}
        & C_2^{k-1} \ar{r}{j} \ar[blue]{u}{d}
        & C_3^{k-1} \ar{r} \ar[blue]{u}{d}
        & 0 \\
        & \vdots \ar[blue]{u}{d}
        & \vdots \ar[blue]{u}{d}
        & \vdots \ar[blue]{u}{d}
        &
        \ar[boundary, from=2-4, to=1-2]{}[']{\delta}
        \ar[boundary, from=3-4, to=2-2]{}[']{\delta}
        \ar[boundary, from=4-4, to=3-2]{}[']{\delta}
    \end{tikzcd}
\] and induce a long exact sequence \[
    \begin{tikzcd}
        \cdots \ar[green!55!black]{r}{\delta} & H^k(C_1) \ar{r}{i^*} & H^k(C_2) \ar{r}{j^*} & H^k(C_3) \ar[green!55!black]{r}{\delta} & H^{k+1}(C_1) \ar{r}{i^*} & \cdots
    \end{tikzcd}
\]

\subsection*{Mayer-Vietoris Sequence\headingnote{Tu, \emph{An Introduction to Manifolds}, \S26}}

We apply the above discussion to the case of de Rham cohomology. From now on, we assume \(H^k(M)\) is the de Rham cohomology \(H_{\dR}^k(M)\).

\begin{theorem}[Mayer-Vietoris]
    Let \(M\) be a manifold and \(U,V\subset M\) be open sets such that \(M = U\cup V\). Then the \emph{Mayer-Vietoris sequence}
    \[
        \tikzset{
            mv zigzag/.style={
                rounded corners=3pt,
                to path={
                let \p1 = (\tikztostart.east),
                    \p2 = (\tikztotarget.west),
                    \n1 = {(\y1+\y2)/2}
                in
                (\p1)
                    -- ({\x1+5mm},\y1)
                    -- ({\x1+5mm},\n1)
                    -- ({\x2-5mm},\n1) \tikztonodes
                    -- ({\x2-5mm},\y2)
                    -- (\p2)
                }
            }
        }
        \begin{tikzcd}[column sep=large, row sep=huge]
            H^{k+1}(M) \ar{r}{i^*} & \cdots \hspace{4em} \\
            H^k(M) \ar{r}{i^*}
            & H^k(U)\oplus H^k(V) \ar{r}{j^*}
            & H^k(U\cap V) \ar[mv zigzag]{ull}[swap]{\delta} \\
            & \hspace{5em}\cdots \ar{r}{j^*} & H^{k-1}(U\cap V) \ar[mv zigzag]{ull}[swap]{\delta}
        \end{tikzcd}
    \] is long exact, where \(i:U\longhookrightarrow M\) is the inclusion and \(j:U\cap V\longhookrightarrow U\oplus V\) is the restriction.
\end{theorem}
\begin{proof}
    Construct the short exact sequence \[
        \begin{tikzcd}
            0 \ar{r} & \Omg^k(M) \ar{r}{i} & \Omg^k(U)\oplus\Omg^k(V) \ar{r}{j} & \Omg^k(U\cap V) \ar{r} & 0 \\
            & \omg \ar[|->]{r}{i} & (\omg|_U, \omg|_V), \\
            & & (\omg_1, \omg_2) \ar[|->]{r}{j} & \omg_1|_{U\cap V} - \omg_2|_{U\cap V}.
        \end{tikzcd},
    \]
\end{proof}

\begin{corollary}\label{cor:mv_euler_char}
    Let \(M\) be a manifold and \(U,V\subset M\) be open sets such that \(M = U\cup V\). Suppose \(\dim H^k(M), \dim H^k(U), \dim H^k(V), \dim H^k(U\cap V)<\infty\), then \[
        \chi(M) - (\chi(U) + \chi(V)) + \chi(U\cap V) = 0.
    \]
\end{corollary}
\begin{proof}
    Write the Mayer-Vietoris sequence as \[
        \begin{tikzcd}
            0 \ar{r}{f_{-1}} & C^0 \ar{r}{f_0} & C^1 \ar{r}{f_1} & C^2 \ar{r}{f_2} & \cdots \ar{r}{f_{n-1}} & C^n \ar{r}{f_n} & 0,
        \end{tikzcd}
    \] then by rank-nullity theorem and exactness, \begin{align*}
        \sum_{k=0}^n (-1)^k\dim C^k
        &= \sum_{k=0}^n (-1)^k(\dim\ker f_k + \dim\im f_k) \\
        &= \sum_{k=0}^n (-1)^k(\dim\im f_{k-1} + \dim\im f_k) \\
        &= \dim\im f_{-1} + \sum_{k=0}^{n-1} (\dim\im f_k - \dim\im f_k) + (-1)^n\dim\im f_n = 0,
    \end{align*} since \(\im f_{-1} = 0 = \im f_n\).
\end{proof}

One favorable case of applying the Mayer-Vietoris sequence is when \(H^k(U)\) and \(H^k(V)\) are simpler, then we can compute \(H^k(M)\).

\tikzcdset{scale cd/.style={every label/.append style={scale=#1},
    cells={nodes={scale=#1}}}}
\begin{example}
    \(H^k(S^n) = \begin{cases}
        \R & k = 0, n,\\
        0 & 1\leq k\leq n-1.
    \end{cases}\).
\end{example}
\begin{proof}
    We will show for \(n\geq 2\), \begin{enumerate}[label=({\itshape \arabic*})]
        \item \(H^1(S^n) = 0\), and
        \item \(H^k(S^n) = H^{k-1}(S^{n-1})\) for \(2\leq k\leq n\).
    \end{enumerate}
    Take the cover \[
        U = S^n\setminus\{\text{north pole}\} \cong \R^n,\quad
        V = S^n\setminus\{\text{south pole}\} \cong \R^n,
    \] by contracting to equator, we see \(U\cap V\simH S^{n-1}\).
    \begin{enumerate}[label=({\itshape \arabic*})]
        \item We have \[
            \begin{tikzcd}[sep=small]
                & 0 \ar{r}
                & H^0(S^n) \ar{r} \ar[equals]{d}
                & H^0(U)\oplus H^0(V) \ar{r} \ar[equals]{d}
                & H^0(U\cap V) \ar{r}{\delta} \ar[equals]{d}
                & H^1(S^n) \ar{r} \ar[equals]{d}
                & H^1(U)\oplus H^1(V) \ar{r} \ar[equals]{d}
                & \cdots \\
                & 0 \ar{r}
                & \R \ar{r}{i}
                & \R\oplus\R \ar{r}{j}
                & \R \ar{r}{\delta}
                & H^1(S^n) \ar{r}
                & 0 \oplus 0
            \end{tikzcd}
        \]
        \begin{itemize}
            \item Injectivity of \(i\) gives \(\R = \im i = \ker j\).
            \item Rank-nullity theorem gives \[
                \dim\im j = \dim(\R\oplus\R) - \dim\ker j = 2 - 1 = 1 = \dim\R,
            \] so \(j\) is surjective.
            \item Surjectivity of \(j\) gives \(\R = \im j = \ker\delta\), so \(\delta = 0\), hence \(H^1(S^n) = 0\).
        \end{itemize}
        \item We have \[
            \begin{tikzcd}[sep=small]
                & \cdots \ar{r}
                & H^{k-1}(U)\oplus H^{k-1}(V) \ar{r} \ar[equals]{d}
                & H^{k-1}(U\cap V) \ar{r}{\delta} \ar[equals]{d}
                & H^k(S^n) \ar{r} \ar[equals]{d}
                & H^k(U)\oplus H^k(V) \ar{r} \ar[equals]{d}
                & \cdots \\
                & 
                & 0\oplus 0 \ar{r}
                & H^{k-1}(S^{n-1}) \ar{r}{\delta}
                & H^k(S^n) \ar{r}
                & 0 \oplus 0
            \end{tikzcd}
        \]
        Injectivity and surjectivity of \(\delta\) gives \(H^k(S^n) = H^{k-1}(S^{n-1})\).
    \end{enumerate}
\end{proof}

\begin{corollary}
    If \(m\neq n\), then \(\R^m\) is not homeomorphic to \(\R^n\).
\end{corollary}
\begin{proof}
    Let \(f:\R^m\longrightarrow\R^n\) be a homeomorphism, then \[
        f:\R^m\setminus\{0\}\simH S^{m-1}\longrightarrow\R^n\setminus\{f(0)\}\simH S^{n-1}
    \] is a homeomorphism, hence \(H^{m-1}(S^{m-1}) = H^{m-1}(S^{n-1})\), which forces \(m = n\).
\end{proof}

% \begin{lemma}
%     If \(\dim H^k(M) < \infty\) for all \(k\) and \[
%         \begin{tikzcd}
%             0 \ar{r} & H^1(M) \ar{r} & H^2(M) \ar{r} & \cdots \ar{r} & H^n(M) \ar{r} & 0
%         \end{tikzcd}
%     \] is exact, then \(\chi(M) = 0\).
% \end{lemma}

\begin{example}
    \(H^k(T^2) = \begin{cases}
        \R & k = 0, \\
        \R^2 & k = 1, \\
        \R & k = 2, \\
        0 & k > 2.
    \end{cases}\)
\end{example}
\begin{proof}
    Since \(T^2\) is connected, we have \(H^0(T^2) = 1\). Since \(\dim T^2 = 2\), we have \(H^k(T^2) = 0\) for \(k > 2\). We will show \begin{enumerate}[label=({\itshape \arabic*})]
        \item \(H^2(T^2) = \R\)
        \item \(H^1(T^2) = \R^2\)
    \end{enumerate}
    Take the cover \[
        U = S^1\times (S^1\setminus\{\text{north pole}\}) \simH S^1,\quad
        V = S^1\times (S^1\setminus\{\text{south pole}\}) \simH S^1,
    \] then \(U\cap V\) consists of two disjoint cylinders, so \(U\cap V\simH S^1\sqcup S^1\). 
    \begin{enumerate}[label=({\itshape \arabic*})]
        \item Since \(H^1(U\cap V) = H^1(S^1\sqcup S^1) = H^1(S^1)\oplus H^1(S^1) = \R\oplus\R\), we have \[
            \begin{tikzcd}[sep=small]
                & \cdots \ar{r}
                & H^1(U)\oplus H^1(V) \ar{r} \ar[equals]{d}
                & H^1(U\cap V) \ar{r}{\delta} \ar[equals]{d}
                & H^2(T^2) \ar{r} \ar[equals]{d}
                & H^2(U)\oplus H^2(V) \ar{r} \ar[equals]{d}
                & \cdots \\
                & 
                & \R\oplus\R \ar{r}{j}
                & \R\oplus\R \ar{r}{\delta}
                & H^2(T^2) \ar{r}
                & 0\oplus 0
            \end{tikzcd}
        \]
    Let \(\omg_U\) and \(\omg_V\) be the generators of \(H^1(U) = \R\) and \(H^1(V) = \R\), respectively, then \begin{align*}
        j: H^1(U)\oplus H^1(V) = \R\oplus\R &\longrightarrow H^1(U\cap V) = \R\oplus\R, \\
        (a\omg_U, b\omg_V) &\longmapsto (a\omg_U,a\omg_V) - (b\omg_U,b\omg_V) = \Big((a-b)\omg_U + (b-a)\omg_V\Big),
    \end{align*} so \(\im j = \R\) is the diagonal line in \(\R^2\). By this and the surjectivity of \(\delta\), we have \[
        \ker\delta = \im j = \R, \quad \im\delta = H^2(T^2).
    \] Hence, rank-nullity theorem gives \[
        2 = \dim \R\oplus\R = \dim\im\delta + \dim\ker\delta = \dim H^2(T^2) + 1,
    \] so \(\dim H^2(T^2) = 1\), hence \(H^2(T^2) = \R\).
    \item We have \[
        \begin{tikzcd}
            0 \ar{r}
            & H^0(T^2) \ar{r} \ar[equals]{d}
            & H^0(U)\oplus H^0(V) \ar{r} \ar[equals]{d}
            & H^0(U\cap V) \ar{r}{\delta} \ar[equals]{d} & \ \\
            & \R
            & \R\oplus\R
            & \R^2 \\
            \ar{r}{\delta} & H^1(T^2) \ar{r}
            & H^1(U)\oplus H^1(V) \ar{r} \ar[equals]{d}
            & H^1(U\cap V) \ar{r}{\delta} \ar[equals]{d} & \  \\
            & & \R\oplus\R
            & \R\oplus\R \\
            \ar{r}{\delta} & H^2(T^2) \ar{r} \ar[equals]{d}
            & H^2(U)\oplus H^2(V) \ar{r} \ar[equals]{d}
            & H^2(U\cap V) \ar{r}{\delta} \ar[equals]{d} & \  \\
            & \R
            & 0
            & 0
        \end{tikzcd}
    \]
    Therefore, by Corollary~\ref{cor:mv_euler_char}, we have \[
        (1 - 2 + 2) + (\dim H^1(T^2) - 2 + 2) + (1 - 0 + 0) = 0,
    \] so \(\dim H^1(T^2) = 2\), hence \(H^1(T^2) = \R^2\).
    \end{enumerate}
\end{proof}

\pagebreak

\section{De Rham Cohomology with Compact Support}

\subsection{Good Covers}

We have seen that if \(U\cap V\) is simple, like \(\emptyset\) or \(\R^n\), which is contractible, then such cover is handy.
\begin{definition}
    We say an open cover \(\{U_\a\}_\a\) is a \emph{good cover} if for any finite subset \(I = \{\a_1,\ldots,\a_k\}\), the intersection \[
        U_I = U_{\a_1}\cap\cdots\cap U_{\a_k}
    \] is either empty or diffeomorphic to \(\R^n\).
\end{definition}


\begin{example}
Consider the following covers for \(S^1\):
\begin{figure}[H]
    \centering
    \tikzset{
        circlebase/.style={draw=black, line width=1pt},
        coverarc/.style={line width=2pt, line cap=butt}
    }

    \begin{subfigure}[t]{0.47\linewidth}
        \centering
        \begin{tikzpicture}[scale=0.88]
            \path[use as bounding box]
                (-2.35,-2.30) rectangle (2.35,2.30);

            \draw[circlebase] (0,0) circle[radius=1.34];
            \node at (0,0) {\(S^1\)};

            % Two open arcs of 200 degrees, with disconnected overlap.
            \draw[coverarc, red]
                (-10:1.67)
                arc[start angle=-10, end angle=190, radius=1.67];

            \draw[coverarc, cyan]
                (170:1.94)
                arc[start angle=170, end angle=370, radius=1.94];
        \end{tikzpicture}
        \caption{}
        \label{fig:s1-bad-cover}
    \end{subfigure}
    \hfill
    \begin{subfigure}[t]{0.47\linewidth}
        \centering
        \begin{tikzpicture}[scale=0.88]
            \path[use as bounding box]
                (-2.35,-2.30) rectangle (2.35,2.30);

            \draw[circlebase] (0,0) circle[radius=1.34];
            \node at (0,0) {\(S^1\)};

            % Three open arcs of 200 degrees, centered 120 degrees apart.
            \draw[coverarc, red]
                (-10:1.67)
                arc[start angle=-10, end angle=190, radius=1.67];

            \draw[coverarc, cyan]
                (110:1.94)
                arc[start angle=110, end angle=310, radius=1.94];

            \draw[coverarc, yellow]
                (230:2.21)
                arc[start angle=230, end angle=430, radius=2.21];
        \end{tikzpicture}
        \caption{}
        \label{fig:s1-good-cover}
    \end{subfigure}
\end{figure}
    The cover in (a) is not good because the intersection of the two arcs is disconnected, hence not \(\R^n\), while the cover in (b) is good because the intersection of any two arcs is still a single open arc, hence diffeomorphic to \(\R^1\).
\end{example}

\begin{proposition}
    If a manifold \(M\) admits a good cover, then \(\dim H^k(M)<\infty\) for all \(k\).
\end{proposition}

\begin{proposition}
    A compact manifold \(M\) always admits a finite good cover.
\end{proposition}

\subsection*{Compact Supports\headingnote{Bott, Tu, \emph{Differential Forms in Algebraic Topology}, \S1.2}}

\begin{definition}
    For \(\omg\in\Omg(M)\), the support of \(\omg\) is defined as \[
        \supp\omg = \overline{\{p\in M:\omg_p\neq 0\}}.
    \] We say \(\omg\) is \emph{compactly supported} if \(\supp\omg\) is compact.
    Denote by \(\Omg_c^k(M)\) the space of compactly supported \(k\)-forms on \(M\).
\end{definition}

With these definitions, we obtain a cochain complex \[
    \begin{tikzcd}
        0 \ar{r} & \Omg_c^0(M) \ar{r}{d} & \Omg_c^1(M) \ar{r}{d} & \cdots \ar{r}{d} & \Omg_c^n(M) \ar{r} & 0,
    \end{tikzcd}
\]

\begin{definition}
    Denote by the vector spaces \[
        Z_c^k(M) := \ker\Big(d:\Omg_c^k(M)\longrightarrow\Omg_c^{k+1}(M)\Big),\quad
        B_c^k(M) := \im\Big(d:\Omg_c^{k-1}(M)\longrightarrow\Omg_c^k(M)\Big),
    \] and the quotient vector space \[
        H_{\dR}^k(M) := Z_c^k(M) \Big/ B_c^k(M)
    \] is the \(k\)-th de Rham cohomology with compact support of \(M\).
\end{definition}

\begin{proposition}\label{prop:Hc0_conn_comp}
    Let \(M\) be a smooth manifold with \(r\) compact connected components, then \(H_c^0(M)\cong \R^r\).
\end{proposition}
\begin{proof}
    Note that \[
        Z_c^0(M) = \{
            f\text{ constant on each connect component of }M\text{ and }\supp(f)\text{ is compact}
        \},
    \] so \(r = \#{\text{compact connected components of }M}\).
\end{proof}

\begin{example}
    \(H_c^k(\R) = \begin{cases}
        0 & k = 0, \\
        \R & k = 1.
    \end{cases}\)
\end{example}

Now we see a big difference between \(H^k(M)\) and \(H_c^k(M)\), for example, \[
    H_c^0(\R) = 0, \text{ but } H^0({\text{point}}) = \R,
\] but \(\R\simH \{\text{point}\}\), so \(H_c^k\) is not a homotopy invariant.

\begin{definition}[\emph{proper maps}]
    A map \(\f:M\longrightarrow N\) is called \emph{proper} if for any compact set \(K\subset N\), the preimage \(f^{-1}(K)\) is compact in \(M\).
\end{definition}
When \(\f\) is proper, we recover the \emph{pullback} \[
    f^*:\Omg_c^k(N)\longrightarrow\Omg_c^k(M).
\]

\begin{theorem}
    Let \(f,g:M\longrightarrow N\) be \emph{properly} homotopic, then \[
        f^* = g^*:H^k(N)\longrightarrow H^k(M).
    \]
\end{theorem}
\begin{proof}
    Exercise.
\end{proof}

\begin{theorem}
    
\end{theorem}

\section{Degree of a Mapping}

Assume \(M\) and \(N\) are connected, oriented manifolds.

\begin{lemma}
    Let \(f\) be a proper map. If \(C\) is closed in \(M\), then \(f(C)\) is closed in \(N\).
\end{lemma}

\begin{proposition}
    Let \(f:M\longrightarrow N\) be a proper map. If \(f\) is not surjective, then \(\deg(f) = 0\).
\end{proposition}
\begin{proof}
    Note that \(f(M)\) is closed subset of \(N\). Hence if \(q\not\in f(M)\), then there exists an open neighborhood \(V_q\) of \(q\) such that \(f^{-1}(V_q) = \emptyset\).
    Take \(\omg\in\Omg_c^n(V_q)\), then \(f^*\omg = 0\), since \[
        0 = \int_M f^*\omg = \deg(f)\int_N\omg,
    \] we have \(\deg(f) = 0\).
\end{proof}

\begin{proposition}
    \(\deg(f)\) is an integer.
\end{proposition}
\begin{proof}
    By Sard's theorem, there exists a \(A\subset N\) such that \(N\setminus A\) is a null set such that for all \(q\in A\) is a regular value of \(f\), i.e.\ \(df|_p\) has \emph{full rank} for all \(p\in f^{-1}(q)\). Since \(M\) and \(N\) have the same dimension, we see \(df_p\) is an isomorphism, hence by inverse function theorem \(f\) is a local diffeomorphism at \(p\).
    
    Since \(f\) is proper, \(f^{-1}(q)\) is compact, then \(f^{-1}(q) = \{p_1,\ldots,p_k\}\).

    Let \[
        \sgm_i = \begin{cases}
            1 & f|_{U_i} \text{ is orientation-preserving},\\
            -1 & f|_{U_i} \text{ is orientation-reversing}.
        \end{cases}
    \]

    We want to show \(\deg(f) = \sum_i \sgm_i\in\Z\).
    \begin{roof}
        Given \(\omg\in\Omg_c^n(V)\), \(f^*\omg\) is supported on \[
            f^{-1}(V) = U_1\sqcup\cdots\sqcup U_k,
        \] then we have \begin{align*}
            \int_M f^*\omg
            = \sum_{i=1}^k \int_{U_i} f^*\omg
            = \sum_{i=1}^k \sgm_i \int_{V} \omg.
        \end{align*}
    \end{roof}
\end{proof}

As we see, if we pertub \(f\) a little bit, the degree should change without jumping, but since it is an integer, it must be constant. This is the following theorem.
\begin{proposition}
    If \(f,g:M\longrightarrow N\) are \emph{properly homotopic}, then \(\deg(f) = \deg(g)\).
\end{proposition}
\begin{proof}
    By homotopic, there exists \(F:M\times I\longrightarrow N\) such that \[
        F(x,0) = f(x),\quad F(x,1) = g(x).
    \]
    Put \(f_t = F(x,t)\), which is a family of proper maps, then \(\deg(f_t)\) is a continuous function into \(\Z\), hence it is constant. Therefore, \[
        \deg(f) = \deg(f_0) = \deg(f_1) = \deg(g).
    \]
\end{proof}

\begin{theorem}[Hopf Degree Theorem]
    Let \(M\) be a \emph{compact} manifold. For \(f,g:M^n \longrightarrow S^n\), \begin{center}
        \(f\) and \(g\) arer homotopic \(\iff \deg(f) = \deg(g)\).
    \end{center}
\end{theorem}
\begin{proof}
    Project.
\end{proof}

\newpage

\begin{proposition}
    Let \(M^n\) be a \(n\)-dimension manifold with boundary and \(X\) a compact \((n-1)\)-dimensional manifold. Let \(f:\pa M\longrightarrow X\), if \(f\) extends to a smoooth map \(g:M\longrightarrow X\), then \(\deg f = 0\).
\end{proposition}
\begin{proof}
    Suppose \(\omg\in\Omg^{n-1}(X)\) with \(\int_X\omg = 1\). By Stokes' theorem, \begin{align*}
        \deg(f) = \int_{\pa M} f^*\omg
        = \int_{\pa M} g^*\omg
        = \int_M dg^*\omg
        = \int_M g^*d\omg
        = \int_M g^*0 = 0
    \end{align*} because \(\omg\) is already a top form on \(X\). Hence \(\deg(f) = 0\).
\end{proof}

We discuss some applications.
\begin{theorem}[Brower's fixed point theorem]
    Let \(f:B^n\longrightarrow B^n\) be a smooth map, then \(f\) admits a fixed point.
\end{theorem}
\begin{proof}
    Suppose \(f(x)\neq x\) for all \(x\in B^n\). For \(x\in B^n\), the ray \[
        (1 - t)f(x) + tx,\quad t\in[0,1]
    \] through \(f(x)\) to \(x\) intersects \(\pa B^n = S^{n-1}\) at a unique point \(\c(x)\). Note that if \(x\in S^{n-1}\), then \(\c(x) = x\) Now we get a smooth map \begin{align*}
        \c: B^n &\longrightarrow S^{n-1} \\
        x &\longmapsto \c(x)
    \end{align*} such that \(\c = \id\) and \(\pa B^n = S^{n-1}\).
    This \(\c\) is an extension of \(\id:\pa B^n\longrightarrow S^{n-1}\), but \(\deg(\id) = 1\), which is a contradiction.
\end{proof}

Now we move on the last essential part of whatever a course of homology should discuss, the Poincaré duality, the symmetry of the cohomology groups.

\section[Poincar\'e Duality]{Poincar\'e Duality\headingnote{Madsen, Tornhave, \emph{From Calculus to Cohomology}, \S13}}

So far we have studied indiviadually the cohomology classes, now we would look at the relationship between all those cohomology in different dimensions. We use the wedge product to relate them.

Let \(M\) be an oriented manifold of \(\dim M = n\), recall that \begin{align*}
    \wedge: H^k(M)\times H_c^l(M)&\longrightarrow H_c^{k+l}(M),\\
     ([\omg],[\eta])&\longmapsto [\omg\wed\eta].
\end{align*}
and if \(M\) is connected, then \begin{align*}
    \int_M\cdot: H_c^n(M) &\longrightarrow \R,\\
    [\omg] &\longmapsto \int_M\omg.
\end{align*}
We have those maps, if we compose them, we get the so called \emph{Poincaré map} \begin{align*}
    P_M^k: H^k(M)\times H_c^{n-k}(M) &\longrightarrow \R,\\
    P_M^k([\omg],[\eta]) &= \int_M \omg\wed\eta,
\end{align*} which is a bilinear map. Recall if we have a bilinear map, then by currying, we can induce a pairing for the first vector space and the dual of the second vector space, i.e.\ \begin{align*}
    \mathcal{P}_M^k: H^n(M) &\longrightarrow \Big(H_c^{n-k}(M)\Big)^*
    = \operatorname{Hom}\Big(H_c^{n-k}(M)\longrightarrow\R\Big) \\
    [\omg] &\longmapsto \Big([\eta]\longmapsto \int_M \omg\wed\eta\Big).
\end{align*}

\begin{theorem}[Poincar\'e duality]
    If \(M\) is an oriented manifold of \(\dim M = n\), then \[
        \mathcal{P}_M^k: H^k(M)\xrightarrow[\hspace{2em}]{\cong} \Big(H_c^{n-k}(M)\Big)^*
    \] is a linear isomorphism. Equivalently, \(P_M^k\) is a \emph{non-degenerate} pairing.
\end{theorem}
\begin{proof}
    Assume \(M\) admits a finite good cover, then we perform induction on the number of charts in the good cover.
\end{proof}

\begin{example}
    If \(M = \R^n\), then \[
        H_{\dR}^k(\R^n) = \begin{cases}
            \R & k = 0,\\
            0 & k > 0,
        \end{cases}
        \qquad
        H_c^k(\R^n) = \begin{cases}
            0 & k < n,\\
            \R & k = n.
        \end{cases}
    \]
\end{example}

\begin{example}
    If \(M = S^n\), then \[
        H_{\dR}^k(S^n) = \begin{cases}
            \R & k = 0, \\
            0 & 0 < k < n, \\
            \R & k = n,
        \end{cases}
        \qquad
        H_c^k(S^n) = \begin{cases}
            \R & k = 0, \\
            0 & 0 < k < n, \\
            \R & k = n.
        \end{cases}
    \]
\end{example}

As a consequence, if \(M\) is a non-compact connected manifold, then by Poincar\'e duality, we have \[
    H^n(M) \cong H_c^0(M)^* = \{0\}^* = 0,
\] for the number of compact connected components.

\begin{example}
    By Poincar\'e duality, we compute \begin{align*}
        H_c^{k+l}(M\times\R^l)
        \underset{PD}{\cong} H^{n-k}(M\times\R^l)^*
        \cong H^{n-k}(M)^*
        \underset{PD}{\cong} H_c^k(M).
    \end{align*}
\end{example}

\begin{example}
    Let \(M\) be a compact oriented manifold of \(\dim M = n\), then \begin{align*}
        b_k = \dim H^k(M) \underset{PD}{=} \dim H_c^{n-k}(M)^* = \dim H_c^{n-k}(M) = b_{n-k}.
    \end{align*}
    Recall that \[
        \chi(M) = \sum_{k=0}^n (-1)^k b_k
    \]
    We have \begin{itemize}
        \item if \(\dim M\) is odd, then \(\chi(M) = 0\),
        \item if \(\dim M \equiv 4m + 2\), then \(\chi(M) \equiv b_{2m+1}\pmod{2}\). Consider the non-degenerate pairing \begin{align*}
            P_M^{2m+1}: H^{2m+1}(M)\times H_c^{2m+1}(M) &\rightarrow \R, \\
            P_M^{2m+1}([\omg],[\eta]) &= \int_M \omg\wed\eta = \int_M (-1)^{2m+1} \eta\wed\omg = -P_M^{2m+1}([\eta],[\omg]),
        \end{align*} then the pairing matrix is anti-symmetric of size \(b_{2m+1}\times b_{2m+1}\), therefore \begin{align*}
            \det(P_M^{2m+1}) = \det((P_M^{2m+1})^t) = \det(-P_M^{2m+1}) = (-1)^{b_{2m+1}}\det(P_M^{2m+1}),
        \end{align*}
        Since \(\det\neq 0\), hence \(b_{2m+1}\) must be even, which implies \(\chi(M)\equiv 0\pmod{2}\).
    \end{itemize}
\end{example}

% \begin{example}
%     By Poincaré duality, we have \[
%         H^1(\R^2\setminus\{0\}) \cong H_c^{1}(\R^2\setminus\{0\})^* \cong H_c^{1}(S^1\times\R)^* \cong H_c^{1}(S^1)^* \cong H^{0}(S^1)^* \cong \R.
%     \]
% \end{example}
